\documentclass{article}
\usepackage{geometry}
\usepackage{concmath}
\usepackage[OT1]{fontenc}
\usepackage{graphicx}
\usepackage{amsfonts}
\usepackage{amsmath,mathtools}
\usepackage{calrsfs}
\usepackage{color}
\usepackage[dvipsnames]{xcolor}
\usepackage[many]{tcolorbox}
\usepackage{setspace}
\usepackage{tkz-tab}    
\usepackage{amssymb}          
\usepackage{multicol}  
\geometry{margin=0.5in}
\definecolor{green}{RGB}{52, 235, 168}
\definecolor{violet}{RGB}{165, 83, 237}
\definecolor{lemon}{RGB}{170, 237, 83}
\definecolor{red}{RGB}{240, 125, 84}
\definecolor{blue}{RGB}{105, 255, 242}

\begin{document} \noindent
\colorbox{green}{230.12}\\
a/ \(\mathcal{B}(e_1)\)\\
\(Mat(v, \mathcal{B}) = (a_{1,1})\)\\
\(v^k = 0\) donc \(a_{1,1}^k\) = 0 \\
donc \(a_{1,1} = 0\)\\
\(\rightarrow v = 0 \)\\
b/ Soit \(\mathcal{B} = (e_1,...,e_p,e_{p+1},...,e_n)\)\\
base de E adaptée à Imv \\
p = dimImv = rg v\\
on a \(1 \leq p < n-1\) car \(v \neq 0\) et v nilpotent donc non bijectif\\
\(u \circ v = v \circ u\) donc Ker v et Imv stables par u \\
Notons \(B = Mat(v, \mathcal{B}), A = Mat(u \mathcal{B})\)
\begin{align*}
    det(u+v) &= \begin{vmatrix}
        A_1 + B+1 & A_2+B_2\\
        0 & A_3
    \end{vmatrix}
    &= det(A_1+B_1)det(A_3)\\
    &= det(A_1)det(A_3)\\
    &= det(u)
\end{align*}
On note \begin{align*}
    v_1 &: Imv \rightarrow Imv \\
    &x \mapsto v(x)
\end{align*}
L'endomorphisme induit par u sur Imv (car Imv stable par v)\\
et \(u_1\) l'endomorphisme induit par u sur Imv car Imv stable par u\\
b/ Soit \(\mathcal{B}_1 = (e_1,...,e_p)\) \\
\(B_1 = Mat(v_1, \mathcal{B_1})\\
A_1 = Mat(u_1, \mathcal{B}_1)\)\begin{align*}
    u_1 \circ v_1 &= u(v(x))\\
    &= v(u(x))\\
    &= v_1 \circ u_1 (x)
\end{align*}
\(\exists k\) tq \(\forall x \in E, v^k(x) = 0\)\\
donc en partie \(\forall x \in E_1, v_1^k(x) = 0\)
Ainsi \(u_1\) et \(v_1\) commutent, \(v_1\) nilpotent\\
Donc par HR forte \begin{align*}
    det(u_1 + v_1) = det(u_1)\\
    ie det(A_1 + B_1) = det(A_1)
\end{align*}
\underline{Formaliser}\\
\(\forall p \in \mathbb{N}^*\) si E est un ev de dim p
u, v \(\in \mathcal{L}(E)\) tq \(u\circ v = v \circ u\), v nilpotent\\
Donc det(u+v) = detu



\end{document}